\section{Open questions (5)}

\begin{enumerate}

\item \textbf{Time-dependent Hamiltonians.}
Does the Trotter24 direct-measurement estimator $\eta^{(24)}$ retain leading-order fidelity to the true infidelity when $H = H(t)$ (Floquet drive, chirped pulses)?
\emph{Basis:} the paper's Eq.~8 leading-order argument assumes a static generator; time-dependent $H(t)$ requires Magnus- or Dyson-series Trotter variants and the estimator's $T_2/T_4$ error-cancellation was not derived for these settings.
\emph{Next steps:} reimplement $T_2/T_4$ for a time-ordered Trotter (e.g.\ midpoint-rule Magnus), test on a linearly-driven Ising chain $H(t)=H_0 + f(t)\,S_x$ at $L{=}8$, and check whether $\eta^{(24)}/\eta_\text{true}$ stays within 10\% across the same $\delta t$ window verified for the static case.

\item \textbf{Order-adaptive Trotter24.}
Can the $T_4$-reference estimator be fused with an adaptive scheme that also adapts the operator \emph{order} (switching between $T_2/T_4/T_6$ at runtime), and does that beat fixed-order Trotter24 on total gate count?
\emph{Basis:} Trotter24 keeps the pair $(2,4)$ fixed; in stiff regions higher-order formulas permit much larger $\delta t$ with fewer total exponentials, but no hybrid-order variant is analyzed.
\emph{Next steps:} prototype an order-adaptive scheduler that picks $(m,n)\in\{(2,4),(4,6),(6,8)\}$ per step based on the local $\eta^{(m,n)}$ estimator; compare against fixed Trotter24 on total-exponentials-per-simulated-time for the mixed-field Ising at $L{=}10$, tolerance $10^{-4}$.

\item \textbf{Chemistry Hamiltonians.}
How does $\eta^{(24)}$ behave on second-quantized chemistry Hamiltonians (H$_2$, LiH, N$_2$ STO-3G) where $A$ and $B$ are chosen by fermionic operator grouping rather than site-local Ising terms?
\emph{Basis:} the paper's mixed-field Ising is nearest-neighbor with short-ranged nested commutators; molecular Hamiltonians (Jordan-Wigner) have highly non-local commutators, and the $T_4$-correction structure may not decouple as cleanly.
\emph{Next steps:} take the OpenFermion H$_2$ Hamiltonian in JW encoding, split $A/B$ by term-anticommutation grouping, run the adaptive Trotter24 loop to a $1$~mHa ground-state target, and report $\eta^{(24)}$ vs true infidelity across $\delta t$; compare adapted step to Childs--Su bound.

\item \textbf{Noise / shot-noise robustness.}
Does the fidelity-overlap estimator $|\langle T_4\psi | T_2\psi\rangle|^2$ retain its precision guarantee under realistic shot noise plus 2-qubit gate error (depolarizing $p\sim 10^{-3}$)?
\emph{Basis:} the paper's analysis is noise-free; on hardware the overlap must be estimated by SWAP/Hadamard tests, both with finite variance and bias from gate error. It is not obvious the estimator remains a good proxy once shot noise dominates over $\eta$ itself.
\emph{Next steps:} simulate the overlap with $10^{4}$--$10^{6}$ shots and a depolarizing channel $p{=}10^{-3}$/2-qubit gate on $L{=}6$ mixed-field Ising; characterize the bias--variance envelope of $\hat{\eta}^{(24)}$ vs noise-free $\eta_\text{true}$ across $\delta t$.

\item \textbf{ML-driven sample-budget allocation.}
Can a small learned surrogate (features: $L$, tolerance, $\|[A,[A,B]]\|$, initial-state entropy, previous $\eta$) predict the per-step shot budget more efficiently than the paper's constant-$C$ rule --- cutting total sample cost by an order of magnitude?
\emph{Basis:} the current rule uses fixed $C{=}0.95$ with no learning across steps; in practice the same $(\psi_0, H)$ produce highly correlated $\eta$ trajectories that a learned allocator could exploit.
\emph{Next steps:} generate a supervised dataset from the reference loop across $\sim$500 random trajectories ($L\in\{4..10\}$, $\varepsilon\in[10^{-4},10^{-2}]$) on the Ising benchmark; train a small regressor predicting shots-per-step for a target $\hat{\eta}$ variance; benchmark total-shots-vs-tolerance-hit-rate against constant-$C$ on a held-out set.

\end{enumerate}
